\documentclass[11pt, UTF8]{ctexart}
\usepackage{amsmath, amsfonts, amssymb} 
\usepackage{geometry}                   
\usepackage{xcolor}                     
\usepackage{enumitem}                   
\usepackage[most]{tcolorbox}            
\usepackage{fancyhdr}                   
\setlength{\headheight}{13.6pt}
\geometry{a4paper, margin=0.8in}        
\pagestyle{fancy}
\fancyhf{}
\lhead{\color{fblue} 3月12日作业 答案} 
\rhead{\color{fblue} \today}
\cfoot{\thepage}

\definecolor{fblue}{RGB}{0, 74, 142}
\definecolor{fred}{RGB}{204, 26, 26}

\newtcolorbox{problem}[2]{
    enhanced,
    colback=white,
    colframe=fblue!80,
    arc=2mm,
    boxrule=0.8pt,
    before upper={\textbf{\color{fblue} 问题 #1 #2}\hspace{0.5em}},
    before skip=1.5em,
    after skip=1.5em,
    left=3mm, right=3mm, top=2mm, bottom=2mm
}

\newcommand{\Prob}{P}
\newcommand{\comp}[1]{\overline{#1}\mskip 3mu}

\begin{document}

\begin{center}
    {\LARGE \bfseries \color{fblue} 3月12日作业 答案}
\end{center}

\begin{problem}{1}{1.3}
（1）设 $A, B, C$ 是三个事件，且 $\Prob(A)=\Prob(B)=\Prob(C)=\frac{1}{4}, \Prob(AB)=\Prob(BC)=0, \Prob(AC)=\frac{1}{8}$，求 $A, B, C$ 至少有一个发生的概率.

（2）已知 $\Prob(A)=\frac{1}{2}, \Prob(B)=\frac{1}{3}, \Prob(C)=\frac{1}{5}, \Prob(AB)=\frac{1}{10}, \Prob(AC)=\frac{1}{15}, \Prob(BC)=\frac{1}{20}, \Prob(ABC)=\frac{1}{30}$，求 $A \cup B, \comp{A}\comp{B}, A \cup B \cup C, \comp{A}\comp{B}\comp{C}, \comp{A}\comp{B}C, \comp{A}\comp{B} \cup C$ 的概率.

（3）已知 $\Prob(A)=\frac{1}{2}$，(i) 若 $A, B$ 互不相容，求 $\Prob(A\comp{B})$，(ii) 若 $\Prob(AB)=\frac{1}{8}$，求 $\Prob(A\comp{B})$.
\end{problem}

\paragraph{\color{fblue}解} 
（1）
\begin{align*}
    \Prob(A \cup B \cup C) &= \Prob(A)+\Prob(B)+\Prob(C)-\Prob(AB)-\Prob(BC)-\Prob(AC)+\Prob(ABC) \\
    &= \frac{5}{8}+\Prob(ABC).
\end{align*}
由 $ABC \subset AB$，已知 $\Prob(AB)=0$，故 $0 \leqslant \Prob(ABC) \leqslant \Prob(AB)=0$，得 $\Prob(ABC)=0$. \\
所求概率为 $\Prob(A \cup B \cup C) = \mathbf{\color{fred}5/8}$.

（2）
\begin{align*}
    \Prob(A \cup B) &= \Prob(A)+\Prob(B)-\Prob(AB) = \frac{1}{2}+\frac{1}{3}-\frac{1}{10} = \mathbf{\color{fred}\frac{11}{15}}.\\
    \Prob(\comp{A}\comp{B}) &= \Prob(\comp{A \cup B}) = 1-\Prob(A \cup B) = \mathbf{\color{fred}\frac{4}{15}}.\\
    \Prob(A \cup B \cup C) &= \Prob(A)+\Prob(B)+\Prob(C)-\Prob(AB)-\Prob(AC)-\Prob(BC)+\Prob(ABC) \\
    &= \frac{1}{2}+\frac{1}{3}+\frac{1}{5}-\frac{1}{10}-\frac{1}{15}-\frac{1}{20}+\frac{1}{30} = \frac{51}{60} = \mathbf{\color{fred}\frac{17}{20}}.\\
    \Prob(\comp{A}\comp{B}\comp{C}) &= \Prob(\comp{A \cup B \cup C}) = 1-\Prob(A \cup B \cup C) = \mathbf{\color{fred}\frac{3}{20}}.\\
        \Prob(\comp{A}\comp{B}C) &= \Prob(\comp{A}\comp{B}(S-\comp{C})) = \Prob(\comp{A}\comp{B}-\comp{A}\comp{B}\comp{C}) = \Prob(\comp{A}\comp{B})-\Prob(\comp{A}\comp{B}\comp{C}) \\
    &= \frac{4}{15}-\frac{3}{20} = \frac{16-9}{60} = \mathbf{\color{fred}\frac{7}{60}}.
\end{align*}
记 $p=\Prob(\comp{A}\comp{B} \cup C)$，由加法公式
\[ p = \Prob(\comp{A}\comp{B})+\Prob(C)-\Prob(\comp{A}\comp{B}C) = \frac{4}{15}+\frac{1}{5}-\frac{7}{60} = \mathbf{\color{fred}\frac{7}{20}}. \]

（3）(i) $\Prob(A\comp{B}) = \Prob(A(S-B)) = \Prob(A-AB) = \Prob(A)-\Prob(AB) = \mathbf{\color{fred}1/2}$. \\
(ii) $\Prob(A\comp{B}) = \Prob(A(S-B)) = \Prob(A-AB) = \Prob(A)-\Prob(AB) = 1/2-1/8 = \mathbf{\color{fred}3/8}$.

\vspace{1.5em}

\begin{problem}{2}{1.6}
在房间里有 10 个人，分别佩戴从 1 号到 10 号的纪念章，任选 3 人记录其纪念章的号码.
\begin{enumerate}
    \item[(1)] 求最小号码为 5 的概率.
    \item[(2)] 求最大号码为 5 的概率.
\end{enumerate}
\end{problem}

\paragraph{\color{fblue}解}
$E$: 在房间里任选 3 人，记录其佩戴的纪念章的号码. 10 人中任选 3 人共有 $\binom{10}{3}=120$ 种选法，此即为样本点的总数. 以 $A$ 记事件“最小的号码为 5”，以 $B$ 记事件“最大的号码为 5”.

（1）因选到的最小号码为 5，则其中一个号码为 5 且其余两个号码都大于 5，它们可从 $6 \sim 10$ 这 5 个数中选取，故 $N(A)=\binom{5}{2}$，从而
\[ \Prob(A) = N(A)/N(S) = \binom{5}{2} \Big/ \binom{10}{3} = \mathbf{\color{fred}\frac{1}{12}}. \]

（2）同理，$N(B)=\binom{4}{2}$，故
\[ \Prob(B) = N(B)/N(S) = \binom{4}{2} \Big/ \binom{10}{3} = \mathbf{\color{fred}\frac{1}{20}}. \]

\vspace{1.5em}

\begin{problem}{3}{1.8}
在 1500 件产品中有 400 件次品，1100 件正品. 任取 200 件.
\begin{enumerate}
    \item[(1)] 求恰有 90 件次品的概率.
    \item[(2)] 求至少有 2 件次品的概率.
\end{enumerate}
\end{problem}

\paragraph{\color{fblue}解}
$E$: 从 1500 件产品中任取 200 件产品. 以 $A$ 表示事件“恰有 90 件次品”，以 $B_i$ 表示事件“恰有 $i$ 件次品”，$i=0,1$，以 $C$ 表示事件“至少有 2 件次品”.

（1）$N(S)=\binom{1500}{200}$，
\[ N(A)=\binom{400}{90}\binom{1100}{200-90}=\binom{400}{90}\binom{1100}{110}, \]
故
\[ \Prob(A) = N(A)/N(S) = \mathbf{\color{fred}\binom{400}{90}\binom{1100}{110} \Big/ \binom{1500}{200}}. \]

（2）$C = S-B_0-B_1$，其中，$B_0, B_1$ 互不相容，所以
\begin{align*}
    \Prob(C) &= \Prob(S-B_0-B_1) = \Prob(S-(B_0 \cup B_1)) \\
    &= 1-\Prob(B_0 \cup B_1) = 1-\Prob(B_0)-\Prob(B_1).
\end{align*}
因
\[ N(B_0)=\binom{1100}{200}, \quad N(B_1)=\binom{400}{1}\binom{1100}{199}, \]
故
\[ \Prob(B_0)=\binom{1100}{200} \Big/ \binom{1500}{200}, \quad \Prob(B_1)=\binom{400}{1}\binom{1100}{199} \Big/ \binom{1500}{200}, \]
因此有
\begin{align*}
    \Prob(C) &= 1-\binom{1100}{200} \Big/ \binom{1500}{200} - \binom{400}{1}\binom{1100}{199} \Big/ \binom{1500}{200} \\
    &= \mathbf{\color{fred}1-\left[ \binom{1100}{200} + \binom{400}{1}\binom{1100}{199} \right] \Big/ \binom{1500}{200}}.
\end{align*}

\end{document} 